\emph{Setting.} Under (B') the joint algebra is the CAR algebra over $\ell^{2}(J) \oplus \ell^{2}(J')$ (Fact \claimref{hawking.imported-quasi-free-states-car-algebra} for its states), its windows are $\mathfrak{A}_{F} \otimes \mathfrak{A}'_{F'} \cong \End(\Lambda^{\bullet}\C^{F}) \otimes \End(\Lambda^{\bullet}\C^{F'})$ (Definition \claimref{obj.hawking.boundary-mode-algebra-windows-boost-dynamics}), with scalars $\mathrm{span}\{1, e_{2}e_{3}\}$, $e_{2}e_{3}$ the boundary modes' complex unit, and the two charge rotations act on a window by conjugation with $e^{\theta e_{2}e_{3}Q_{F}}$ and $e^{\theta' e_{2}e_{3}Q'_{F'}}$, $Q_{F} = \sum_{j \in F}N_{j}$, $Q'_{F'} = \sum_{j \in F'}N'_{j}$. Representing the superselection of Proposition \claimref{dark-matter.superselected-hence-dark-exactly-stable} by these invariants is the hypothesis (B'); the stored asymmetry is the eigenvalue of such a charge (Proposition \claimref{dark-matter.stored-asymmetry-charge-not-number}(i)).

(i) An element of $\End(\Lambda^{\bullet}\C^{F})$ is invariant under conjugation by $e^{\theta e_{2}e_{3}Q_{F}}$ for all $\theta$ if and only if it commutes with $Q_{F}$, that is, if and only if it preserves each eigenspace $\Lambda^{q}\C^{F}$; so the invariant subalgebra is $\bigoplus_{q}\End(\Lambda^{q}\C^{F})$, a direct sum of full matrix algebras, whose centre is spanned by the eigenprojections $P_{q}$ of $Q_{F}$. The same holds for the mirror, and the observable window algebra, the tensor product of the two invariant subalgebras, is $\bigoplus_{q, q'}\End(\Lambda^{q}\C^{F} \otimes \Lambda^{q'}\C^{F'})$ with centre spanned by the $P_{q} \otimes P'_{q'}$. Two window states supported in different joint sectors are separated by the central projection of one of them, so they are disjoint. The check realises a window of two modes per tower in $\Cl(4,4)$, confirms the CAR relations, confirms that the charge projections are idempotent, sum to one, have ranks $4$, $8$, $4$ of $16$, and commute with every bilinear $a^{*}_{i}a_{j}$ of either tower, and computes the dimension of the algebra those bilinears generate as $36 = (1 + 4 + 1)^{2}$, the dimension of $\bigoplus_{q,q'}\End$ for sectors of dimensions $1, 2, 1$ in each tower; it also exhibits two states of different charge separated by a central projection.

(ii) For one-particle Hamiltonians $h$, $h'$ and inverse temperatures $\beta$, $\beta'$ the product of Gibbs states has density $\rho = \prod_{j \in F}(1 + (x_{j} - 1)N_{j}) \prod_{j \in F'}(1 + (x'_{j} - 1)N'_{j})$ up to normalisation, $x_{j} = e^{-\beta\omega_{j}}$, $x'_{j} = e^{-\beta'\omega'_{j}}$, all positive; its inverse is the same product with every $x$ replaced by $1/x$, so $\rho$ is invertible and the state is faithful. A faithful state gives every nonzero projection positive weight, in particular every $P_{q} \otimes P'_{q'}$, so its central support is $1$. Two states with the same central support are not separated by any central projection: if $Z$ is central with $\omega(Z) = 1$ then $\omega(1 - Z) = 0$, so $1 - Z = 0$ by faithfulness, and $\omega'(Z) = 1$ as well. On a finite-dimensional algebra this is quasi-equivalence. So within every joint sector the ratio $\xi$ is a parameter of the state and not of the sector. The check confirms $\rho\rho^{-1} = 1$ exactly, that all nine sector weights are positive and sum to one, and that the mirror's occupation is $x'/(1 + x')$ and differs from ours exactly when $x' \neq x$, for $x = \tfrac13$ and $x' \in \{\tfrac13, \tfrac{1}{50}, \tfrac27\}$.

(iii) The I-state at $\beta$ is the quasi-free state with diagonal symbol $n_{j} = (e^{\beta\omega_{j}} + 1)^{-1}$ (Theorem \claimref{hawking.i-state}), gauge-invariant. The Powers-Stormer criterion (Fact \claimref{synthesis.imported-powers-stormer-takesaki-disjointness}(a)) states that two gauge-invariant quasi-free states with symbols $A$, $B$ are quasi-equivalent if and only if $A^{1/2} - B^{1/2}$ and $(1 - A)^{1/2} - (1 - B)^{1/2}$ are Hilbert-Schmidt, and that such states are factor states, so that they are disjoint whenever they are not quasi-equivalent. For diagonal symbols the sum of the two squared Hilbert-Schmidt norms is $\sum_{j}\bigl[(\sqrt{n_{j}} - \sqrt{n'_{j}})^{2} + (\sqrt{1 - n_{j}} - \sqrt{1 - n'_{j}})^{2}\bigr] = 2\sum_{j}(1 - a_{j})$, which gives the stated dichotomy. For a mode of frequency $\omega_{j} > 0$, $n_{j} \neq n'_{j}$ when $\beta \neq \beta'$, and $a_{j} < 1$: squaring, $a_{j}^{2} < 1$ is equivalent to $4n n'(1 - n)(1 - n') < (n + n' - 2nn')^{2}$, and directly, $a_{j}$ is the Bhattacharyya coefficient of two distinct probability vectors $(n, 1 - n)$ and $(n', 1 - n')$, which equals $1$ only when they coincide (Cauchy-Schwarz with equality only for proportional vectors). The bound $1 - a_{j} \le |n_{j} - n'_{j}|$ follows from $(\sqrt{p} - \sqrt{q})^{2} \le |\sqrt{p} - \sqrt{q}|(\sqrt{p} + \sqrt{q}) = |p - q|$ applied to both terms. Now $1 - a$ is a continuous, strictly positive function of $\omega$ on $(0, \infty)$ at fixed $\beta \neq \beta'$, so it is bounded below by some $c > 0$ on any compact subset; infinitely many frequencies in such a subset make $\sum_{j}(1 - a_{j}) \ge \sum c = \infty$, and the states are disjoint. If instead $\sum_{j}(n_{j} + n'_{j}) < \infty$, then $\sum_{j}(1 - a_{j}) \le \sum_{j}|n_{j} - n'_{j}| < \infty$ and they are quasi-equivalent. Disjoint factor states of the full algebra are separated by a central projection of the von Neumann algebra of their direct sum, so in the first case $\beta$ is a label in the disjointness sense. The same conclusion for KMS factor states whose induced von Neumann algebra is of type III is Takesaki's theorem (Fact \claimref{synthesis.imported-powers-stormer-takesaki-disjointness}(b)), the standard result behind (iii). The comparison with (i) needs care: there the charge is a label because it is the eigenvalue of a central element of a finite window algebra. On the full algebra $Q = \sum_{j}N_{j}$ is not an element of the observable C*-algebra, and a thermal I-state over infinitely many modes has no finite total charge; what can label such a state is a charge density, and only in the same disjointness sense (by the criterion above with $n_{j}$ shifted by a chemical potential). So the stored charge $Q = \epsilon N$ is a finite sector label of the deposited state exactly when the deposited mode set is finite, the same finiteness the reading below invokes for the temperature; with infinitely many modes both are labels in the disjointness sense, and a scope to labels reaches both. The check verifies, in exact rational arithmetic over a grid of seven Gibbs parameters $x = e^{-\beta\omega}$ in each argument, that $a^{2} < 1$ exactly when $x \neq x'$ (equality exactly on the diagonal), through the rational form above, and that $1 - a \le |n - n'|$, through the rational inequality $4nn'(1 - n)(1 - n') \ge \bigl((1 - |n - n'|)^{2} - nn' - (1 - n)(1 - n')\bigr)^{2}$ whenever the bracket is positive.

\emph{The reading.} Proposition \claimref{synthesis.mirror-asymmetry-route-one} offers, without proof, a scoping under which P2 equates the superselected charges of the descent, which the stored asymmetry is and the temperature ratio is not. By (i) and (ii) that distinction is a representation-theoretic one on every window: the stored charge selects a central projection, the temperature does not. By (iii) it survives the inductive limit exactly when the sum converges; when infinitely many boundary modes share a frequency band it does not, the temperature is then a label of the same kind as the charge, and a scoping to labels equates the two towers' temperatures, $\xi = 1$, which Proposition \claimref{synthesis.cold-mirror-needs-non-equilibrium-descent} reads, under its hypothesis (M), as $\Delta\Neff = 7.45$, excluded. The series counts about $\pi N^{2}/\ln 2$ boundary modes for the cosmological horizon (at observation level, in the Standard-Model paper's largeness observation, not used as a premise here), a finite number, which places the series in the first case if that count bounds the deposited sector's modes; the Definition's countable $J$ does not decide it. The ceiling coupling of Theorem \claimref{dark-matter.mirror-tower-copy-one}(iv) is not a charge, so the scoping does not reach it; it reaches the asymmetry clause (v).
